\section{Prior effective sample size: derivation and computation}\label{sec:B}
\subsection{Normal reference experiment}\label{sec:B.1}
The reference variance $\sigma_2^2$ below is held fixed. In numerical calculations it is replaced by a fixed plug-in estimate $\widehat\sigma_2^2$, distinguished from the variance sampled in the final outcome model. Consider $m$ independent reference observations with common mean $\mu$ and known variance $\sigma_2^2$. Under the flat prior $p(\mu)\propto1$, the posterior kernel is
\begin{equation}
p(\mu\mid Y_{1:m}^{\mathrm{ref}})\propto
\exp\left\{-\frac{1}{2\sigma_2^2}\sum_{j=1}^m(Y_j^{\mathrm{ref}}-\mu)^2\right\}
\propto\exp\left\{-\frac{m}{2\sigma_2^2}(\mu-\bar Y^{\mathrm{ref}})^2\right\}.
\end{equation}
The second proportionality uses $\sum_j(Y_j^{\mathrm{ref}}-\mu)^2=\sum_j(Y_j^{\mathrm{ref}}-\bar Y^{\mathrm{ref}})^2+m(\mu-\bar Y^{\mathrm{ref}})^2$. Normalization gives
\begin{equation}
\mu\mid Y_{1:m}^{\mathrm{ref}}\sim\N(\bar Y^{\mathrm{ref}},\sigma_2^2/m).
\end{equation}
The matching quantity is this \emph{posterior variance of a mean}. The reference count is subsequently allowed to be positive and noninteger; it is an information-equivalence index rather than an actual sample size requirement.

\subsection{Historical-data-informed distribution and historical variance}\label{sec:B.2}
Trial standardization averages over $\mathbf{X}_1$ and $\mathbf{X}_2$, with $N\coloneqq n_1+n_2$. Condition throughout on $\mathbf{X}_3,\mathbf{X}_1,\mathbf{X}_2$ and all fixed prior settings. The preliminary fit uses the likelihood of $\mathbf{Y}_3$ alone. Its posterior distribution for $f$ is combined with an independent prior for $g$. This constructs an informative prior for the trial-control mean before using trial-control outcomes. The final two-arm analysis additionally incorporates trial-control outcomes.

For $\mu_f\coloneqq N^{-1}\sum_{s=1}^2\sum_{i=1}^{n_s} f(x_{s,i})$,
\begin{equation}\label{eq:Vf}
V_f\coloneqq \Var(\mu_f\mid \mathbf{Y}_3,\mathbf{X}_3,\mathbf{X}_1,\mathbf{X}_2)
=\frac{1}{N^2}\sum_{s,t=1}^2\sum_{i=1}^{n_s}\sum_{j=1}^{n_t}\Cov\{f(x_{s,i}),f(x_{t,j})\mid \mathbf{Y}_3,\mathbf{X}_3,\mathbf{X}_1,\mathbf{X}_2\}.
\end{equation}
Given posterior prediction draws $f^{(d)}(x_{s,i})$, compute $\mu_f^{(d)}\coloneqq N^{-1}\sum_{s=1}^2\sum_{i=1}^{n_s} f^{(d)}(x_{s,i})$ and take their sample variance. This integrates historical-tree and leaf uncertainty through posterior sampling. The posterior sample variance incorporates tree uncertainty and cross-profile covariance; averaging pointwise variances would omit the latter.

Trial covariates define the estimand and may enter the rule for candidate cutpoints. Section B.9 specifies the differing cutpoint rules used for the preliminary fit, prior-tree simulation and final fit.

\subsection{Discrepancy variance: node weights, indicators and trees}\label{sec:B.3}
For the collection $\mathcal T^g\coloneqq (\mathcal T_1^g,\ldots,\mathcal T_{H_g}^g)$ of discrepancy trees, define
\begin{equation}\label{eq:weights}
a_{j\ell}\coloneqq \frac1N\sum_{s=1}^2\sum_{i=1}^{n_s}\ind\{x_{s,i}\in\ell\text{ of tree }j\},\qquad
\mu_g\coloneqq \frac1N\sum_{s=1}^2\sum_{i=1}^{n_s}g(x_{s,i})=\sum_{j,\ell}a_{j\ell}\theta^g_{j\ell}.
\end{equation}
The equality follows by regrouping the sum of tree predictions according to which profiles fall in each node. Within each tree, $a_{j\ell}\geq0$ and $\sum_\ell a_{j\ell}=1$. These weights come from trial profiles passed through \emph{prior} discrepancy trees in the calibration calculation, not from the posterior trees of the historical fit.

Conditional on the trees and indicators, independent zero-mean terminal-node priors imply
\begin{equation}\label{eq:conditionalVg}
\begin{aligned}
\E(\mu_g\mid\mathcal T^g,\mathbf z,w,\tau_0^2)&=0,\\
\Var(\mu_g\mid\mathcal T^g,\mathbf z,w,\tau_0^2)
&=\sum_{j,\ell}a_{j\ell}^2\{z_{j\ell}\tau_0^2+(1-z_{j\ell})\tau_1^2\}.
\end{aligned}
\end{equation}
The law of total variance first averages over the indicators. Because every conditional mean is zero, its between-mean variance term vanishes. Using $\E(z_{j\ell}\mid w)=w$,
\begin{equation}
\Var(\mu_g\mid\mathcal T^g,w,\tau_0^2)
=\{w\tau_0^2+(1-w)\tau_1^2\}\sum_{j,\ell}a_{j\ell}^2.
\end{equation}
Apply the same law over the prior trees; their conditional mean is also zero. Therefore
\begin{equation}\label{eq:Vg}
\boxed{V_g(w,\tau_0^2)
\coloneqq \Var(\mu_g\mid \mathbf{X}_3,\mathbf{X}_1,\mathbf{X}_2,w,\tau_0^2)
=\{w\tau_0^2+(1-w)\tau_1^2\}\E_{\mathcal T^g}\!\left[\sum_{j,\ell}a_{j\ell}^2\right].}
\end{equation}
$V_g$ is the variance of $\mu_g$ under the discrepancy prior after integrating terminal-node, indicator and tree uncertainty at fixed $w,\tau_0^2$. This step is exact under the stated independent zero-mean priors; Monte Carlo approximation enters when evaluating the tree expectation.

\subsection{From a matched count to the ESS definition}\label{sec:B.4}
Independence of the preliminary historical posterior and the discrepancy prior gives
\begin{equation}
\Var(\mu\mid \mathbf{Y}_3,\mathbf{X}_3,\mathbf{X}_1,\mathbf{X}_2,w,\tau_0^2)=V_f+V_g(w,\tau_0^2).
\end{equation}
Equating this to the normal-reference posterior variance and solving produces the conditional matched count
\begin{equation}\label{eq:m}
m(w,\tau_0^2)=\frac{\sigma_2^2}{V_f+V_g(w,\tau_0^2)}.
\end{equation}
At this point $m$ is defined for one fixed spike variance. A scale parameter $s_0^2$ specifies a distribution of such variances. The proposed ESS is the mean of their matched counts:
\begin{equation}\label{eq:ESSoriginal}
\begin{aligned}
\ESS(w,s_0^2)&\coloneqq \E_{\tau_0^2}\!\left[
\frac{\sigma_2^2}{\Var(\mu\mid \mathbf{Y}_3,\mathbf{X}_3,\mathbf{X}_1,\mathbf{X}_2,w,\tau_0^2)}\right]\\
&=\E_{\tau_0^2}\!\left[\frac{\sigma_2^2}{V_f+V_g(w,\tau_0^2)}\right].
\end{aligned}
\end{equation}
The expectation over indicators and trees is \emph{inside the variance}, before taking its reciprocal; the subsequent expectation over the spike variance is outside the ratio. Interchanging these operations defines a different summary.

In particular, when the fully scale-marginalized variance is finite, its variance-matched ESS would instead be $\sigma_2^2/[V_f+\E_{\tau_0^2}\{V_g(w,\tau_0^2)\}]$. By convexity of the reciprocal, it is no larger than \eqref{eq:ESSoriginal}. Neither identity provides the predictive-additivity guarantee of the expected local-information-ratio construction \citep{neuenschwander2020}.

For prior ESS calibration set $w=1$ and use $\tau_0^2\sim\operatorname{scaled\text{-}Inv}\chi^2(\nu_0,s_0^2)$ without truncation. Then
\begin{equation}\label{eq:ESStau}
\ESS_{\tau_0}(s_0^2)\coloneqq\ESS(1,s_0^2)=\E_{\tau_0^2}\!\left[
\frac{\sigma_2^2}{V_f+\tau_0^2\E_{\mathcal T^g}[\sum_{j,\ell}a_{j\ell}^2]}\right].
\end{equation}
This is the reduction of the original definition, with $V_g(1,\tau_0^2)$ written explicitly. All-spike calibration is a reference for choosing a prior scale; it does not replace posterior indicator sampling in the two-arm fit. It is the narrow-component reference, not an exact upper envelope over mixture weights at scale draws exceeding the fixed slab variance.

\subsection{Proof of the prior ESS curve theorem and ceiling}\label{sec:B.5}
Let $U\sim\chi^2_{\nu_0}$ and couple all scale values by $\tau_0^2=\nu_0s_0^2/U$. The tree-weight expectation is positive because every finite tree has at least one occupied node, and is at most $H_g$ because $\sum_\ell a_{j\ell}^2\leq1$. For each $U>0$, the integrand in \eqref{eq:ESStau} is continuous, strictly decreasing and strictly convex as a function of $s_0^2>0$. It is bounded by $\sigma_2^2/V_f$.

Dominated convergence proves continuity and gives
\begin{equation}
\lim_{s_0^2\downarrow0}\ESS_{\tau_0}(s_0^2)=\frac{\sigma_2^2}{V_f},\qquad
\lim_{s_0^2\to\infty}\ESS_{\tau_0}(s_0^2)=0.
\end{equation}
Pointwise strict inequalities yield strict monotonicity and strict convexity after taking expectations. The intermediate value theorem and strict monotonicity establish existence and uniqueness for a target in $(0,\sigma_2^2/V_f)$. This proves main-text Theorem 1 without requiring a finite mean for $\tau_0^2$.

The limit corresponds to suppressing prior discrepancy variation while retaining historical posterior uncertainty about $f$. Hence the ceiling is an information limit for the particular standardized mean, not the historical head count. The assumptions $V_f>0$ and fixed reference variance matter. A finite upper truncation of the scale prior has a different large-scale limit; a grid capped at 10 likewise need not span every theoretical target.

\subsection{Monte Carlo evaluation and implemented prior-scale selection}\label{sec:B.6}
Simulate prior discrepancy trees using the conditioned-on candidate rules. For independent identically distributed trees, the average of their squared-weight sums, multiplied by $H_g$, estimates $\E_{\mathcal T^g}[\sum_{j,\ell}a_{j\ell}^2]$. Estimate $V_f$ from historical posterior averages. For $B$ independent chi-square draws $U_b$, $b=1,\ldots,B$, evaluate
\begin{equation}\label{eq:MCess}
\widehat\ESS_{\tau_0}(s_0^2)\coloneqq \frac1B\sum_{b=1}^B
\frac{\widehat\sigma_2^2}
{\widehat V_f+(\nu_0s_0^2/U_b)\widehat\E_{\mathcal T^g}[\sum_{j,\ell}a_{j\ell}^2]}.
\end{equation}
Reusing the $U_b$ values over a grid makes the Monte Carlo curve monotone. This computation simulates prior trees and scale variables; the preliminary $f$ predictions are posterior draws conditional on the one historical dataset being calibrated. It does not repeatedly generate entire outcome datasets to define the curve.

\noindent\textbf{Algorithm B1: full-variance ESS evaluation and target selection.}
\begin{enumerate}
\item Specify the target profiles, preliminary historical outcome likelihood, candidate splitting rules and fixed reference residual variance estimate.
\item Fit the historical model, form each posterior target-average prediction and estimate $V_f$ from its variance.
\item Simulate discrepancy trees from the specified prior and estimate the expected sum of squared target-profile weights.
\item Draw $U_b\sim\chi^2_{\nu_0}$ and calculate \eqref{eq:MCess} over candidate scales.
\item For an attainable target, solve the monotone equation or select a grid approximation; record requests outside the curve's achievable range separately.
\end{enumerate}

The \emph{implemented prior-scale selector} modifies steps 2 and 5. It divides the historical mean draws into 20 consecutive blocks, calculates each block variance, and rescales those variances so that their arithmetic mean equals the variance across all draws. Write $J=20$ for the number of blocks. Let the rescaled values be $\widehat V_{f,k}$. It evaluates a curve $\widehat\ESS_k$ for each block by substituting that variance in \eqref{eq:MCess}; the reference residual variance is held constant across blocks. The block-based central curve is their arithmetic mean. For a capped target $m_*\coloneqq \min(m_{\mathrm{requested}},\widehat\sigma_2^2/\widehat V_f)$, the selector chooses the smallest supplied scale satisfying
\begin{equation}\label{eq:blockrule}
\frac1J\sum_{k=1}^J\ind\{\widehat\ESS_k(s_0^2)\leq m_*\}\geq0.95.
\end{equation}
If no scale satisfies the condition, it chooses the largest grid value. If the requested target reaches or exceeds the full-variance ceiling, the current rule overrides this search and chooses the smallest grid value. At the boundary, the selected scale need not attain the requested ESS or enforce an upper bound on the block-based curve.

Even after rescaling the block variances to the whole variance, their reciprocal average differs:
\begin{equation}
\lim_{s_0^2\downarrow0}\frac1J\sum_k\widehat\ESS_k(s_0^2)
=\frac1J\sum_k\frac{\widehat\sigma_2^2}{\widehat V_{f,k}}
\ \geq\ \frac{\widehat\sigma_2^2}{\widehat V_f}.
\end{equation}
Equality requires equal positive block variances. The full-variance ceiling is therefore not generally the zero-scale limit of the block-based central curve. Blocks of an autocorrelated MCMC chain also do not furnish a formal 95\% posterior or repeated-sampling guarantee. Figure S1 shows the curves actually used; Table S11 additionally evaluates the full-variance expression. Equation \eqref{eq:MCess} estimates the theoretical prior ESS curve; the reported fits use the block-based selection procedure.

\begin{figure}[p]
\centering
\includegraphics[width=.94\linewidth,height=7.8in,keepaspectratio]{figureS1.pdf}
\caption{Prior ESS curves used for scale selection. Simulation panels use Sc1, replicate 1; the application panel uses PFS. Solid curves are block-based central ESS summaries over the variance-scale grid, with a MAP overlay only where a matching prior ESS calculation is available. Dotted horizontal lines are the full-variance variance ceilings. The lrcBART central curve averages block-specific reciprocal variances, so its small-scale limit need not equal the dotted ceiling. The application curves differ in $H_f$. The MAP and lrcBART methods have different prior constructions and ESS implementations; the shared horizontal scale denotes each method's supplied variance-scale parameter, not identical prior distributions. Each curve uses one dataset; simulation curves are not averaged over replicates.}
\end{figure}

\subsection{Pointwise and subgroup targets}\label{sec:B.7}
For a single profile $x$, one node of each discrepancy tree has weight one. Thus
\begin{equation}
V_g(w,\tau_0^2;x)=H_g\{w\tau_0^2+(1-w)\tau_1^2\},\qquad
\ESS_{\tau_0}(s_0^2;x)=\E_{\tau_0^2}\!\left[\frac{\sigma_2^2}{V_f(x)+H_g\tau_0^2}\right].
\end{equation}
For a subgroup $G$, let $\mathbf{X}_G$ contain its trial profiles, retaining multiplicities, and let $N_G\coloneqq |\mathbf{X}_G|$. Replace the target by $\mu_G\coloneqq N_G^{-1}\sum_{x\in \mathbf{X}_G}\{f(x)+g(x)\}$. Calculate $V_f(G)$ from posterior averages within $G$, and use $a_{j\ell}(G)\coloneqq N_G^{-1}\sum_{x\in \mathbf{X}_G}\ind(x\in\ell)$ in the tree expectation. This gives
\begin{equation}
\ESS_{\tau_0}(s_0^2;G)=\E_{\tau_0^2}\!\left[
\frac{\sigma_2^2}{V_f(G)+\tau_0^2\E_{\mathcal T^g}[\sum_{j,\ell}a_{j\ell}(G)^2]}\right].
\end{equation}
Subgroup comparisons hold the historical fit and globally selected scale fixed and use all observed covariates within each group. The target is the subgroup mean rather than the average of individual-profile ESS values. A singleton group reduces to the pointwise formula.

\subsection{Interpretation and nonadditivity}\label{sec:B.8}
An ESS of 20 for a subgroup means that, averaged over the spike-variance reference distribution, the full historical-data-informed prior provides the precision of approximately 20 independent normal-reference observations for that subgroup's mean. At each fixed scale, it is the posterior variance under the flat-prior reference experiment that is matched. It does not imply that the fully marginalized prior variance is $\sigma_2^2/20$.

For a partition into groups with weights $p_G\coloneqq N_G/N$, $\mu=\sum_Gp_G\mu_G$. At fixed $w,\tau_0^2$ under the same historical-information distribution,
\begin{equation}\label{eq:nonadd}
\Var(\mu)=\sum_Gp_G^2\Var(\mu_G)+\sum_{G\ne H}p_Gp_H\Cov(\mu_G,\mu_H).
\end{equation}
The same fitted functions inform several groups, so their covariance generally is nonzero. Matching this overall variance and then averaging the conditional counts does not equal summing the subgroup ESS values. Nor do individual-profile ESS values generally sum or average to overall ESS. There is no general additive identity for these information summaries.

\subsection{Single-arm and survival references; differences from the fitting prior}\label{sec:B.9}
In a single-arm design, the historical likelihood is the only likelihood for $f$ and the discrepancy remains prior-driven. The reference variance assumes $\sigma_2^2=\sigma_3^2$. Censoring enters the preliminary historical likelihood for survival outcomes, so $V_f$ reflects uncertainty in that fit. However, the normal reference experiment remains uncensored on the log-time scale. Its ESS cannot be read as the number of patients in a reference trial with a particular censoring rate without an additional calibration.

The mathematical construction conditions on an explicit covariate-dependent tree prior. The preliminary historical fits choose cutpoints from historical covariates, while calibration discrepancy-tree simulation uses combined historical and trial covariates. Final two-arm control fitting uses the control covariates available to that fit. The prior-tree simulator limits depth to 12 and does not impose the final two-arm fitting requirement of at least five trial controls per discrepancy leaf. Consequently the expected squared-weight sum is evaluated under the calibration prior, which is not uniformly identical to the final fitted tree support. Thus the prior ESS calculation and final fitting model use different tree supports.

Prior ESS calibration uses an untruncated spike-variance distribution. The default two-arm sampler uses its upper-truncated counterpart; single-arm and application fits instead fix the spike variance at $\min(s_0^2,\tau_1^2/2)$. An all-spike ESS averaged over the calibration law is therefore a design reference for those analyses, not the exact variance-matched ESS of their fixed-scale fitting prior. The primary PFS subgroup table uses the specified calibration definition at the globally selected scale, and is labeled accordingly. Finally, post-fit ESS based on updated discrepancy parameters is a separate summary and is not reported as prior ESS.
